\section{Symbolic Planning}
\label{sec:symbolic_planning}

Symbolic planning models sequential decision making using an explicit representation of states, actions, and goals. The environment is described using a first-order language consisting of predicate symbols and function symbols. Predicate symbols define Boolean properties of objects, while function symbols define numeric quantities.

A fluent is a state variable whose value may change over time. Boolean fluents are obtained by grounding predicate symbols with objects in the domain, and numeric fluents are obtained from function symbols. Let $\Fluents$ denote the set of all fluents. A symbolic state $\SymState \in \SymStates$ is defined as a  over $\Fluents$, assigning each Boolean fluent a value in $\{0,1\}$ and each numeric fluent a value in $\Real$. In propositional domains, a symbolic state can be equivalently represented as the set of grounded predicates that are true, with all others assumed false under the closed-world assumption.

The dynamics of the symbolic model are defined by a finite set of operators $\Operators$. Each operator $\Operator \in \Operators$ is defined as $\Operator = \langle \Pre_{\Operator}, \Eff_{\Operator} \rangle$, where $\Pre_{\Operator}$ is a logical condition over fluents specifying when the operator is applicable, and $\Eff_{\Operator}$ specifies how the state changes after execution. An operator is applicable in a state $\SymState$ if $\SymState \models \Pre_{\Operator}$. Effects modify Boolean fluents through add and delete operations and update numeric fluents through assignments or arithmetic expressions. Operators are typically defined as parameterized schemas and instantiated over objects to produce grounded operators.

A planning task is defined as $\PlanningTask = \langle \Fluents, \Operators, \SymState_0, \Goal \rangle$, where $\SymState_0 \in \SymStates$ is the initial state and $\Goal$ is a logical condition over fluents. A plan is a finite sequence of operators $\Plan = [\Operator_1, \dots, \Operator_n]$. Executing a plan from the initial state produces a sequence of symbolic states obtained by sequentially applying applicable operators. A plan is valid if execution from $\SymState_0$ results in a state that satisfies $\Goal$. Classical symbolic planning assumes deterministic transitions, where the effects of an operator uniquely determine the resulting state.

Planning domains and problems are commonly specified using the Planning Domain Definition Language (PDDL), which provides a standard representation for predicates, functions, operators, and planning tasks~\cite{goel2022rapidlearn}.

% In the case of everything being propositional, the symbolic state space $\SymStates$ is finite, and the planning problem can be solved using search algorithms such as breadth-first search, depth-first search, or heuristic search. In domains with numeric fluents, the state space may be infinite, and specialized algorithms are required to handle numeric constraints and optimization objectives. Symbolic planning has been successfully applied in various domains, including robotics, logistics, and automated reasoning, where it enables the generation of interpretable plans that can be executed by agents to achieve complex goals.      
